%% Strona redakcyjna tomu \dash{} werso strony tytulowej tomu.
%%
%% Rozni sie od strony redakcyjnej wydania handlowego dokladnie w tych trzech
%% miejscach, w ktorych rozni sie tom: nazywa tom, drukuje ISBN, gdy zostal
%% nadany, i nie twierdzi, ze ksiazka ukazuje sie w "dwoch formatach", co
%% przestalo byc prawda w chwili, gdy powstal trzeci.

\thispagestyle{empty}
\vspace*{\fill}

{\small\color{mgrey}\raggedright\setlength{\parindent}{0pt}%
\emph{\mfaseriestitle} \dash{} \lblVolume~\mfavolnumeral,
\emph{\mfavoltitlehere}. Wydanie \bookedition.

\medskip
Copyright \textcopyright{} \bookyear{} \mfabookauthor. Wszelkie prawa
zastrzeżone poza tym, co objęte licencjami poniżej.

\medskip
Tekst na licencji CC BY-NC-SA 4.0. Kod \dash{} każdy skrypt w katalogu
\texttt{code/} i każdy listing, który ta książka z niego drukuje \dash{} na
licencji MIT.

\medskip
\ifx\mfavolisbn\relax
Wydanie własne autora. Dla tego tomu nie nadano numeru ISBN, więc można go
czytać i kopiować na powyższych licencjach, ale nie może jeszcze zostać
skatalogowany przez bibliotekę.
\else
Wydanie własne autora. ISBN \mfavolisbn.
\fi

\medskip
\lblAiAssist

\medskip
Liczby przypięte \pinnedon. To inna data niż rok praw autorskich powyżej
i znaczy co innego.

\medskip
Metoda nauczania programowanego pochodzi od K.~A. Strouda. Ta książka nie jest
z nim powiązana ani z niego wyprowadzona.
\par}

\cleardoublepage
